% =====================================================================
%  Byapar — Complete Shop Owner User Guide
%  Compile with XeLaTeX (preferred) or LuaLaTeX.  See README.md.
%  Uses only standard TeX Live / Overleaf packages so it builds anywhere.
% =====================================================================
\documentclass[11pt,a4paper,openany]{report}

\usepackage{fontspec}                 % XeLaTeX / LuaLaTeX
\usepackage[a4paper,top=2.5cm,bottom=2.4cm,left=2.2cm,right=2.2cm,headheight=16pt]{geometry}
\usepackage{xcolor}
\usepackage{graphicx}
\usepackage[most]{tcolorbox}
\usepackage{enumitem}
\usepackage{booktabs}
\usepackage{array}
\usepackage{tabularx}
\usepackage{longtable}
\usepackage{titlesec}
\usepackage{fancyhdr}
\usepackage{tikz}
\usetikzlibrary{calc,arrows.meta}
\usepackage[hidelinks]{hyperref}
\usepackage{bookmark}

% ---------- Brand palette (from the real app: --primary hsl(158,64%,30%)) ----------
\definecolor{brandgreen}{HTML}{1C7D59}
\definecolor{brandgreendark}{HTML}{145C41}
\definecolor{navy}{HTML}{13293D}
\definecolor{panel}{HTML}{F2F6F4}
\definecolor{panelborder}{HTML}{D4E3DB}
\definecolor{amber}{HTML}{B45309}
\definecolor{amberbg}{HTML}{FDF4E3}
\definecolor{redink}{HTML}{B3261E}
\definecolor{redbg}{HTML}{FCEDEC}
\definecolor{inkgray}{HTML}{4B5563}
\definecolor{hairline}{HTML}{E3E8E6}

% ---------- Fonts: keep the default Latin Modern so it always compiles.
%            A dedicated font only for the Rupee glyph (DejaVu ships with TeX Live). ----------
\newfontfamily\rupeefont{DejaVu Sans}
\newcommand{\rupee}{{\rupeefont \symbol{"20B9}}}  % ₹

\hypersetup{
  pdftitle={Byapar — Complete Shop Owner User Guide},
  pdfauthor={Byapar},
  colorlinks=true, linkcolor=brandgreendark, urlcolor=brandgreendark,
  bookmarksnumbered=true, bookmarksopen=true,
}

% ---------- Headings ----------
\titleformat{\chapter}[display]
  {\sffamily\bfseries\color{navy}}
  {\raggedright\color{brandgreen}\large CHAPTER \thechapter}
  {6pt}
  {\Huge\raggedright}
\titlespacing*{\chapter}{0pt}{6pt}{18pt}

\titleformat{\section}
  {\sffamily\bfseries\large\color{navy}}{\thesection}{0.6em}{}
\titleformat{\subsection}
  {\sffamily\bfseries\normalsize\color{brandgreendark}}{\thesubsection}{0.6em}{}
\setlength{\parindent}{0pt}
\setlength{\parskip}{6pt}

% ---------- Header / footer ----------
\newcommand{\wordmark}{%
  \begin{tikzpicture}[baseline=-0.55ex]
    \fill[brandgreen,rounded corners=2.2pt] (0,-0.9ex) rectangle (2.3ex,1.5ex);
    \draw[white,line width=0.5pt] (0.45ex,0.1ex) -- (0.45ex,1.0ex) -- (1.15ex,1.55ex) -- (1.85ex,1.0ex) -- (1.85ex,0.1ex) -- cycle;
    \node[anchor=west] at (2.7ex,0.3ex) {\sffamily\bfseries\color{navy}Byapar};
  \end{tikzpicture}%
}
\pagestyle{fancy}
\fancyhf{}
\fancyhead[L]{\wordmark}
\fancyhead[R]{\sffamily\footnotesize\color{inkgray}Shop Owner User Guide}
\fancyfoot[C]{\sffamily\footnotesize\color{inkgray}\thepage}
\renewcommand{\headrulewidth}{0.4pt}
\renewcommand{\headrule}{\hbox to\headwidth{\color{hairline}\leaders\hrule height \headrulewidth\hfill}}
\fancypagestyle{plain}{\fancyhf{}\fancyfoot[C]{\sffamily\footnotesize\color{inkgray}\thepage}\renewcommand{\headrulewidth}{0pt}}

% ---------- Callout boxes ----------
\newtcolorbox{notebox}{enhanced,breakable,colback=panel,colframe=brandgreen,
  boxrule=0pt,leftrule=3pt,arc=3pt,left=10pt,right=10pt,top=8pt,bottom=8pt,
  fonttitle=\sffamily\bfseries\color{brandgreendark},title={Good to know}}
\newtcolorbox{tipbox}{enhanced,breakable,colback=panel,colframe=brandgreendark,
  boxrule=0pt,leftrule=3pt,arc=3pt,left=10pt,right=10pt,top=8pt,bottom=8pt,
  fonttitle=\sffamily\bfseries\color{brandgreendark},title={Tip}}
\newtcolorbox{warnbox}{enhanced,breakable,colback=amberbg,colframe=amber,
  boxrule=0pt,leftrule=3pt,arc=3pt,left=10pt,right=10pt,top=8pt,bottom=8pt,
  fonttitle=\sffamily\bfseries\color{amber},title={Important}}
\newtcolorbox{examplebox}{enhanced,breakable,colback=white,colframe=panelborder,
  boxrule=0.8pt,arc=3pt,left=10pt,right=10pt,top=8pt,bottom=8pt,
  fonttitle=\sffamily\bfseries\color{navy},title={Worked example}}
\newtcolorbox{faqbox}[1]{enhanced,breakable,colback=panel,colframe=panelborder,
  boxrule=0.8pt,arc=3pt,left=10pt,right=10pt,top=7pt,bottom=8pt,
  fonttitle=\sffamily\bfseries\color{navy},title={#1}}

% Steps list
\newlist{steps}{enumerate}{1}
\setlist[steps]{leftmargin=2.1em,label=\protect\stepnum{\arabic*},itemsep=5pt,topsep=4pt}
\newcommand{\stepnum}[1]{\raisebox{-1.5pt}{\colorbox{brandgreen}{\makebox[1.35em][c]{\sffamily\bfseries\footnotesize\textcolor{white}{#1}}}}}

% A soft "what you'll see" chip
\newcommand{\ui}[1]{\colorbox{panel}{\rule[-3pt]{0pt}{11pt}\texttt{\textcolor{navy}{#1}}}}

% Screenshot: uses the real image if present in screenshots/, else a tidy placeholder.
\newcommand{\shot}[3]{% {number}{filename}{caption}
  \begin{center}
  \IfFileExists{screenshots/#2}{%
    \fcolorbox{hairline}{white}{\includegraphics[width=0.93\linewidth]{screenshots/#2}}%
  }{%
    \begin{tcolorbox}[width=0.93\linewidth,colback=panel,colframe=panelborder,
       boxrule=0.8pt,arc=4pt,sharp corners=south,halign=center,valign=center,height=4.6cm,
       enhanced,borderline={0.8pt}{0pt}{panelborder,dashed}]
       {\sffamily\bfseries\color{navy}Screenshot to add}\\[3pt]
       {\sffamily\footnotesize\color{inkgray}place \texttt{screenshots/#2} here}
    \end{tcolorbox}%
  }\\[3pt]
  {\sffamily\footnotesize\color{inkgray}\textbf{Figure #1.}~#3}
  \end{center}
}

\begin{document}
\sffamily\color{black}

% =====================================================================
%  COVER
% =====================================================================
\begin{titlepage}
\thispagestyle{empty}
\begin{tikzpicture}[remember picture,overlay]
  \fill[brandgreen] (current page.north west) rectangle ($(current page.north east)+(0,-6.2cm)$);
  \fill[brandgreendark] ($(current page.north west)+(0,-6.2cm)$) rectangle ($(current page.north east)+(0,-6.6cm)$);
  % wordmark on the band
  \node[anchor=north west] at ($(current page.north west)+(2.2cm,-1.6cm)$)
    {\begin{tikzpicture}
       \fill[white,rounded corners=4pt] (0,0) rectangle (1.0,1.0);
       \draw[brandgreen,line width=1.1pt] (0.2,0.18) -- (0.2,0.62) -- (0.5,0.86) -- (0.8,0.62) -- (0.8,0.18) -- cycle;
       \node[anchor=west,white] at (1.25,0.5) {\sffamily\bfseries\Huge Byapar};
     \end{tikzpicture}};
  \node[anchor=west,white] at ($(current page.north west)+(2.25cm,-3.3cm)$)
    {\sffamily\large Business \& Shop Management};
\end{tikzpicture}

\vspace*{7.4cm}
{\sffamily\bfseries\color{navy}\fontsize{27}{31}\selectfont Complete Shop Owner\\[2pt] User Guide}\\[10pt]
{\sffamily\color{brandgreendark}\Large Manage Your Business with Byapar}\\[18pt]
{\sffamily\color{inkgray}\large Your step-by-step handbook for billing, purchases,\\
inventory, payments, parties, and business reports.}

\vfill
\begin{tcolorbox}[colback=panel,colframe=panelborder,boxrule=0.8pt,arc=4pt,left=14pt,right=14pt,top=10pt,bottom=10pt]
{\sffamily\bfseries\color{navy}Who this guide is for}\\[2pt]
{\color{inkgray}Kirana and grocery shops, retail counters, wholesale traders and small
distributors — and the staff who do the billing. No accounting background needed.}
\end{tcolorbox}
\vspace{6pt}
{\sffamily\footnotesize\color{inkgray}Document version 1.0\hfill Covers the Byapar shop application}
\end{titlepage}

% =====================================================================
%  CONTENTS
% =====================================================================
\pagenumbering{roman}
{\sffamily\bfseries\color{navy}\Huge Contents}\par\vspace{10pt}
{\color{brandgreen}\rule{\linewidth}{1.2pt}}\par\vspace{4pt}
\makeatletter
\@starttoc{toc}
\makeatother
\clearpage
\pagenumbering{arabic}

% =====================================================================
\chapter{Welcome to Byapar}
% =====================================================================
Byapar is one simple place to run the money side of your shop. Instead of paper
bills in a drawer, a stock register in one notebook and a ``who owes me'' list in
another, everything lives together and stays in agreement.

\section{What Byapar does for you}
With Byapar you can:
\begin{itemize}[leftmargin=1.3em,itemsep=2pt]
  \item Make bills for customers in seconds, and print or share them.
  \item Record what you buy from suppliers.
  \item Know your stock at any moment — what you have, and what is running low.
  \item Keep a single list of \textbf{parties} (the people and businesses you deal with).
  \item Track money in and money out, and who still owes whom.
  \item See clear reports of how the shop is doing — on screen, or as PDF and Excel.
  \item Scan a photo of a paper bill and let Byapar read it for you.
\end{itemize}

\section{The problems it solves}
\begin{itemize}[leftmargin=1.3em,itemsep=2pt]
  \item \textbf{Scattered paper bills} — every sale and purchase is saved and searchable.
  \item \textbf{Manual stock registers} — stock moves by itself as you bill and buy.
  \item \textbf{``Udhaar'' notebooks} — each party's balance is kept for you, always current.
  \item \textbf{Month-end guessing} — reports show sales, profit and dues at a glance.
\end{itemize}

\section{How it all works together}
The modules are not separate silos. One action updates everything it should:

\begin{center}
\begin{tikzpicture}[font=\sffamily\footnotesize,node distance=6mm,
   box/.style={draw=panelborder,fill=panel,rounded corners=3pt,inner sep=5pt,minimum width=2.6cm,align=center},
   gbox/.style={draw=brandgreen,fill=white,rounded corners=3pt,inner sep=5pt,minimum width=2.6cm,align=center,text=navy},
   arr/.style={-{Latex[length=2mm]},brandgreen,thick}]
  \node[gbox](sale){\textbf{Sale}};
  \node[box,right=1.4cm of sale](stockd){Stock goes\\\textbf{down}};
  \node[box,below=7mm of stockd](recv){Customer\\\textbf{owes you}};
  \node[gbox,below=1.4cm of sale](purch){\textbf{Purchase}};
  \node[box,right=1.4cm of purch](stocku){Stock goes\\\textbf{up}};
  \node[box,below=7mm of stocku](pay){You\\\textbf{owe supplier}};
  \draw[arr](sale)--(stockd);
  \draw[arr](sale)|-(recv);
  \draw[arr](purch)--(stocku);
  \draw[arr](purch)|-(pay);
  \node[box,right=1.4cm of stockd,yshift=-7mm,minimum width=2.9cm](rep){\textbf{Reports}\\read all of it};
  \draw[arr](recv)-|(rep);
  \draw[arr](pay)-|(rep);
\end{tikzpicture}
\end{center}

\begin{notebox}
You never do this maths by hand. You record the sale or purchase; Byapar updates
stock, the party's balance and the reports for you.
\end{notebox}

\section{Why small shops like it}
It is built for a busy counter: big buttons, works on a phone, keeps the exact
rupee figures so there are no end-of-day surprises, and speaks in plain words
rather than accounting jargon.

% =====================================================================
\chapter{Logging In and the Dashboard}
% =====================================================================
\section{Signing in}
Open the web address your business was given and you will see the \textbf{Byapar}
sign-in screen.
\begin{steps}
  \item Enter the \textbf{email} your business set up for you.
  \item Enter your \textbf{password}.
  \item Press \ui{Sign in}.
\end{steps}
\shot{1}{01-login.png}{The Byapar sign-in screen. Use the email and password your business gave you.}

\begin{warnbox}
Use the account your business created for you. The shop owner account can do
everything; a staff (cashier) account may be limited to billing. If you see
``Cannot reach the server'', check your internet and try again.
\end{warnbox}

\section{The dashboard (Home)}
After signing in you land on \textbf{Home}. It shows the day at a glance — recent
sales, money due, and shortcuts to the things you do most.
\shot{2}{02-dashboard.png}{The Home dashboard: your day at a glance.}

\section{Finding your way around}
On the left is the \textbf{sidebar}, grouped so related tasks sit together:
\begin{center}
\small
\begin{tabularx}{\linewidth}{>{\bfseries\color{navy}}l X}
\toprule
Overview & Home, Quick Billing, Bill Import \\
Sales & Sales, Invoices, Sales Returns \\
Purchase & Purchases, Purchase Bills, Purchase Returns, Payables \\
Stock \& Inventory & Stock Dashboard, Products / Items, Stock In, Stock Out, Stock Adjustments, Low Stock, Offers \\
Parties & All Parties, Add Party, Credit Book \\
Money & Money Received, Money Paid, Expenses, Cash \& Bank \\
Reports & Sales, Purchases, Stock, Parties, Expenses, Profit \& Loss, Billers' Day (and GST, if you use GST) \\
\bottomrule
\end{tabularx}
\end{center}
\shot{3}{03-sidebar.png}{The sidebar, grouped by the work you do.}

\begin{tipbox}
Press \ui{Ctrl\,+\,K} (or \ui{Cmd\,+\,K} on a Mac) anywhere to open the command
menu and jump straight to a page or action.
\end{tipbox}

\section{Profile and signing out}
Your name sits at the top right. Open it to see your profile or to \textbf{sign out}.

\section{What to check every morning}
\begin{itemize}[leftmargin=1.3em,itemsep=2pt]
  \item Yesterday's sales on the dashboard.
  \item \textbf{Low Stock} — what to re-order.
  \item \textbf{Parties} report / Credit Book — who still owes you.
\end{itemize}

% =====================================================================
\chapter{First-Time Setup}
% =====================================================================
Do these once, in this order, and daily billing becomes effortless. A new shop
already comes with sensible defaults, so some steps may only need a quick check.

\begin{steps}
  \item \textbf{Business profile \& settings} — open \textbf{Settings} and check your shop name, address and phone.
  \item \textbf{Tax (GST)} — turn GST on only if your shop is GST-registered. If you are not, leave it off and your bills stay simple, with no tax added.
  \item \textbf{Categories} — group your products (e.g.\ Grocery, Personal Care, Cold Drinks).
  \item \textbf{Units} — the units you sell in (piece, kg, litre, packet).
  \item \textbf{Products} — add the items you sell (see Chapter 4).
  \item \textbf{Store} — most shops have one store, and Byapar uses it automatically. Add more only if you keep separate godowns.
  \item \textbf{Opening stock} — enter what you already have on the shelf (see Chapter 4).
  \item \textbf{Parties} — add the customers and suppliers you deal with (see Chapter 5).
  \item \textbf{Opening balances} — if a party already owes you (or you owe them), record that starting amount.
  \item \textbf{Start billing} — you are ready.
\end{steps}

\begin{notebox}
You do not have to finish setup before you start. Add a product or a party the
moment you need it, right from the billing screen.
\end{notebox}

% =====================================================================
\chapter{Products and Inventory}
% =====================================================================
``Stock'' simply means how much of each product you have in the shop right now.
Byapar keeps this count for you.

\section{Adding a product}
Open \textbf{Stock \& Inventory $\rightarrow$ Products / Items} and choose
\ui{Add product}. The fields are:
\begin{center}
\small
\begin{tabularx}{\linewidth}{>{\bfseries\color{navy}}l X}
\toprule
Name & What you call the item, e.g.\ \emph{Sunflower Oil 1L}. \\
Item code / SKU & A short code for the item (optional). \\
Barcode & The barcode number printed on the pack, if any (optional). \\
Category & The group it belongs to. \\
Unit & How you sell it — piece, kg, litre\ldots \\
GST \& HSN/SAC & Only shown if you use GST: the tax rate and tax code. \\
Purchase price & What you normally pay for it. \\
Selling price & What you normally sell it for. \\
Minimum stock level & When stock falls to this, it appears in \textbf{Low Stock}. \\
Description & Any extra note (optional). \\
\bottomrule
\end{tabularx}
\end{center}
\shot{4}{04-add-product.png}{Adding a product with its purchase and selling price.}

\begin{notebox}
The \textbf{selling price} you set here fills in automatically when you add the
item to a bill, so you do not have to remember the rate every time. You can still
change it on any bill.
\end{notebox}

\section{Combo products}
A \textbf{combo} is one item on the bill made of several products (for example a
``Diwali Pack''). When you sell the combo, each product inside it leaves stock on
its own. A combo cannot itself be purchased or counted — only its contents are.

\section{Opening stock}
To tell Byapar what you already have, record it as a \textbf{purchase} (or a
\textbf{Stock In}) at the price it cost you. From then on, stock moves by itself.

\section{Seeing and moving stock}
\begin{itemize}[leftmargin=1.3em,itemsep=2pt]
  \item \textbf{Stock Dashboard} — the live count and value of everything.
  \item \textbf{Stock In / Stock Out} — add or remove stock directly.
  \item \textbf{Stock Adjustments} — fix a count after breakage, loss or a physical check.
  \item \textbf{Low Stock} — items at or below their minimum level, to re-order.
\end{itemize}
\shot{5}{05-stock-dashboard.png}{The Stock Dashboard: what you hold and what it is worth.}

\begin{examplebox}
\textbf{Sunflower Oil 1L} — purchase price \rupee115, selling price \rupee140,
opening stock \textbf{50}.
\begin{itemize}[leftmargin=1.3em,itemsep=1pt]
  \item You \textbf{buy 20} more from a supplier $\rightarrow$ stock becomes \textbf{70}.
  \item You \textbf{sell 5} to a customer $\rightarrow$ stock becomes \textbf{65}.
\end{itemize}
You never change the number by hand — recording the purchase and the sale does it.
\end{examplebox}

\begin{warnbox}
\textbf{Stock transfers} between two stores are not part of the app today. If you
keep more than one godown, move stock with Stock In / Stock Out at each store.
\end{warnbox}

% =====================================================================
\chapter{Unified Parties}
% =====================================================================
This is one of the most important ideas in Byapar, and it makes life simpler.

\section{What a ``party'' is}
A \textbf{party} is any person or business you deal with — a customer, a supplier,
or someone who is \emph{both}. There are no separate ``customer'' and ``supplier''
lists to keep in step. There is \textbf{one party}, and Byapar remembers both
sides of your dealings with them.

\section{Adding and opening a party}
Open \textbf{Parties $\rightarrow$ Add Party}, enter the name (phone and the
rest are optional), and save. Open any party from \textbf{All Parties} to see their
profile, balances and full history.
\shot{6}{06-add-party.png}{Adding a party — one record for a customer, a supplier, or both.}
\shot{7}{07-party-list.png}{All Parties: everyone you deal with, with what they owe you and what you owe them.}

\section{Two balances, kept apart}
A party can owe you money \emph{and} you can owe them money at the same time.
Byapar keeps these two figures \textbf{separate} and never quietly cancels one
against the other — because the two debts are real and settled separately.

\begin{examplebox}
\textbf{XYZ Traders} — you deal with them both ways:
\begin{center}
\small
\begin{tabular}{@{}l r@{}}
You sell them goods & \rupee8,000 \\
They pay you & $-$\,\rupee2,000 \\
\textbf{Owes You} & \textbf{\rupee6,000} \\[4pt]
You buy goods from them & \rupee20,000 \\
You pay them & $-$\,\rupee5,000 \\
\textbf{You Owe} & \textbf{\rupee15,000} \\
\end{tabular}
\end{center}
Both numbers stand on their own: \textbf{Owes You \rupee6,000} and
\textbf{You Owe \rupee15,000}. Byapar does not net them to ``\rupee9,000''.
\end{examplebox}

\section{The party statement}
Every party has a full \textbf{statement} — one running account of sales,
purchases, receipts and payments, with separate running ``owes you'' and ``you
owe'' totals. You can download it as PDF or Excel (see Chapter 13).

% =====================================================================
\chapter{Quick Billing}
% =====================================================================
Quick Billing is the fastest way to make a sale at the counter. Open
\textbf{Overview $\rightarrow$ Quick Billing}.

\section{Step by step}
\begin{steps}
  \item \textbf{Choose the party} (the customer). New face? Tap \ui{New party? Add them first} and add them in seconds.
  \item \textbf{Add items.} Pick a product; its \textbf{price fills in automatically} from the selling price. Set the \textbf{quantity}. Add as many items as you like.
  \item Change the \textbf{rate} or add a \textbf{discount} on any line if needed. (To give an item free, tick the free option — the stock still goes out.)
  \item Choose how it is paid: \ui{Full} (paid now), \ui{Part} (some now, rest on credit) or \ui{On credit} (pay later), and the account the money goes \ui{Into}.
  \item Press \ui{Save \& complete} to finish the bill, or \ui{Save draft} to keep it for later.
\end{steps}
\shot{8}{08-quick-billing.png}{Quick Billing: choose the party, add items, take payment, complete.}

\begin{notebox}
The \textbf{total, tax and change} are worked out by Byapar when you save, and
shown on the finished invoice straight after — so the figure is always exact.
\end{notebox}

\begin{examplebox}
\textbf{Rahul Store} buys:
\begin{center}
\small
\begin{tabular}{@{}l r r r@{}}
\toprule
Item & Qty & Rate & Amount \\
\midrule
Sunflower Oil 1L & 2 & \rupee140 & \rupee280 \\
Dove Soap 100g   & 3 & \rupee55  & \rupee165 \\
Maggi Noodles 70g& 5 & \rupee14  & \rupee70  \\
\midrule
\multicolumn{3}{r}{\textbf{Items total}} & \textbf{\rupee515} \\
\bottomrule
\end{tabular}
\end{center}
If your shop does \textbf{not} use GST, the bill total is \textbf{\rupee515}.
If you \textbf{do} use GST, Byapar adds the tax on top when you save. Either way,
the final figure appears on the invoice.
\end{examplebox}

\section{What happens when the sale completes}
When you press \ui{Save \& complete}: the stock of each item \textbf{goes down},
the \textbf{invoice} is created, any amount not paid becomes what the customer
\textbf{owes you}, and the money paid is recorded into your chosen account.

% =====================================================================
\chapter{Sales and Invoices}
% =====================================================================
\section{Making a sale the full way}
Beyond Quick Billing, \textbf{Sales $\rightarrow$ New sale} gives you the full
form (dates, payment terms, per-line discounts). A sale is first saved as a
\textbf{draft} — nothing moves yet.
\begin{steps}
  \item Pick the \textbf{party}, store and date.
  \item Add items — the \textbf{price fills in} from each product; set quantity and any discount. A live \textbf{``lines add up to''} figure shows as you type.
  \item Press \ui{Save draft} to keep it, or \ui{Save \& complete} to post it now (owner).
\end{steps}
\shot{9}{09-new-sale.png}{New sale: a running total shows as you add items; one button saves, one completes.}

\begin{notebox}
\textbf{Draft vs completed.} A draft changes nothing and can be edited or deleted.
\textbf{Completing} (posting) a sale moves the stock and the money, and the invoice
becomes final. To change a completed invoice, use a \textbf{sales return}.
\end{notebox}

\section{Invoices}
\textbf{Sales $\rightarrow$ Invoices} lists every completed sale with its
\textbf{payment status} and \textbf{outstanding} amount. Open any invoice to view,
\textbf{print} or \textbf{download} it.

\section{Sales returns}
When a customer returns goods, record a \textbf{sales return}. The returned items
come \textbf{back into stock} and the customer's balance is reduced accordingly.

% =====================================================================
\chapter{Purchases and Purchase Bills}
% =====================================================================
\section{Recording a purchase}
Open \textbf{Purchase $\rightarrow$ Purchases $\rightarrow$ New purchase}.
\begin{steps}
  \item Choose the \textbf{party} (the supplier) and enter the \textbf{supplier's bill number} and date.
  \item Add items with their \textbf{purchase rate}, quantity and any discount.
  \item Choose what you paid now (or leave it for credit).
  \item \ui{Save draft}, or \ui{Save \& complete} to record it (owner).
\end{steps}
\shot{10}{10-new-purchase.png}{New purchase: add items at their purchase rate and record what you paid.}

\section{What a completed purchase does}
It \textbf{adds the items into stock} (at their cost), and any amount you did not
pay becomes what \textbf{you owe} that supplier. Your average cost per item is kept
up to date automatically.

\section{Purchase bills, returns and payables}
\begin{itemize}[leftmargin=1.3em,itemsep=2pt]
  \item \textbf{Purchase Bills} — the list of everything you have bought, with paid and due amounts.
  \item \textbf{Purchase Returns} — goods sent back to a supplier; stock goes down and your due to them is reduced.
  \item \textbf{Payables} — how much you owe each party, and \emph{which bills are due} (see Chapter 10).
\end{itemize}

% =====================================================================
\chapter{Money Received and Money Paid}
% =====================================================================
These two screens record cash and bank movements with your parties. They are kept
separate because one is money coming \textbf{in} and the other money going
\textbf{out}.

\section{Money Received (from customers)}
Open \textbf{Money $\rightarrow$ Money Received} and \ui{Record money received}.
\begin{steps}
  \item Choose the \textbf{party}.
  \item Enter the \textbf{amount} received. Byapar \textbf{applies it to their open bills automatically}, oldest first — you type the figure once.
  \item Choose \textbf{how} it was paid (Cash, UPI, Bank, Cheque) and the account it goes \textbf{into}.
  \item Add a \textbf{reference} or \textbf{note} if you wish.
  \item Press \ui{Record payment} to post it, or \ui{Save draft}.
\end{steps}
\shot{11}{11-money-received.png}{Money Received: the amount fills the open bills on its own; anything extra is kept as an advance.}

\begin{notebox}
Anything left over after the open bills are covered is kept as an \textbf{advance}
for that party, and used against their next bill. You can also split the amount
across bills by hand if you prefer.
\end{notebox}

\section{Money Paid (to suppliers)}
\textbf{Money $\rightarrow$ Money Paid} works the same way, for money you pay
a supplier. Choose the party, enter the amount (it fills their open bills), pick
the method and account, and \ui{Record payment}.

\begin{examplebox}
A customer owes \rupee235 on one invoice. You receive \rupee235 and press
\ui{Record payment}: the invoice is marked paid and the customer's ``owes you''
drops to \rupee0. Receive \rupee300 instead, and \rupee235 settles the invoice
while \rupee65 is kept as an advance.
\end{examplebox}

% =====================================================================
\chapter{Credit Book and Payables}
% =====================================================================
\section{The Credit Book}
\textbf{Parties $\rightarrow$ Credit Book} is your ``udhaar'' register, done
for you. It shows who has an open balance so you can follow up — clearly split
between \textbf{receivable} (money customers owe you) and \textbf{payable} (money
you owe suppliers).

\begin{notebox}
\textbf{Receivable} = money a customer still has to pay \emph{you}.\\
\textbf{Payable} = money you still have to pay a \emph{supplier}.
\end{notebox}

\section{Payables}
\textbf{Purchase $\rightarrow$ Payables} is the matching view for money you
owe: how much you owe each party, which \textbf{bills are due}, the oldest bill,
and a \ui{Pay} button to settle it.
\shot{12}{12-credit-book.png}{Credit Book / Payables: who owes whom, and what is due.}

% =====================================================================
\chapter{AI Bill Import}
% =====================================================================
\textbf{Overview $\rightarrow$ Bill Import} lets you photograph a bill and have
Byapar read it, so you do not have to type it in.

\section{How it works}
\begin{steps}
  \item Press \ui{Upload a bill} and take a photo, or choose a PDF. Printed and neat handwritten bills both work.
  \item Byapar \textbf{reads} the bill — the party, the items, quantities, rates, discount and totals.
  \item \textbf{Check what was read} on the review screen. Fix anything misread, match each line to one of your products, and pick the supplier (or customer).
  \item Press \textbf{confirm} to save it as a \textbf{purchase} (a bill you bought) or a \textbf{sale}.
\end{steps}
\shot{13}{13-bill-import.png}{Bill Import: upload a bill, review what was read, then confirm.}

\begin{warnbox}
\textbf{Always check before you confirm.} The reader is very good, but it can
misread a handwritten figure. Nothing is saved until you press confirm — so review
the party, the items and the totals first. Byapar also \textbf{warns you} if a bill
looks like one you have already recorded, so you do not enter it twice.
\end{warnbox}

\begin{notebox}
\textbf{Discounts are read too.} If the bill has a discount, Byapar shows it in a
\textbf{Discount on the bill} box and spreads it across the items when you save,
so the recorded total matches the bill. A bill with no discount is unaffected.
\end{notebox}

\begin{warnbox}
If you see ``Automatic bill reading is not set up on this server'', the reading
service has not been switched on for your shop yet. You can still upload bills to
keep the photo on file, and enter them with the normal Purchase and Sales forms in
the meantime.
\end{warnbox}

% =====================================================================
\chapter{Expenses, Cash and Bank}
% =====================================================================
\section{Recording an expense}
\textbf{Money $\rightarrow$ Expenses} is for money the shop spends that is
\emph{not} a payment to a supplier — rent, electricity, tea, transport. Choose the
category, enter the amount, and the account it was paid from.

\begin{notebox}
\textbf{Expense vs supplier payment.} Paying your \emph{supplier} for goods is
\textbf{Money Paid} (it reduces what you owe them). Paying \emph{rent} or the
\emph{electricity bill} is an \textbf{Expense}. Keeping them apart keeps your
profit figure honest.
\end{notebox}

\section{Cash \& Bank}
\textbf{Money $\rightarrow$ Cash \& Bank} shows the balance of your cash box and
bank/UPI accounts, and lets you move money between them. Every receipt, payment
and expense updates these balances.

% =====================================================================
\chapter{Reports and Business Performance}
% =====================================================================
Open \textbf{Reports} and use the tabs across the top: \textbf{Sales},
\textbf{Purchases}, \textbf{Stock}, \textbf{Parties}, \textbf{Expenses},
\textbf{Profit \& Loss} (and \textbf{GST}, if you use GST). Most reports let you
pick a \textbf{date range} and \ui{Print}, \ui{Download} or \ui{Share}.

\section{What each report tells you}
\begin{center}
\small
\begin{tabularx}{\linewidth}{>{\bfseries\color{navy}}l X}
\toprule
Sales & What you sold, over any period; returns and net sales. \\
Purchases & What you bought, what you paid, what you still owe. \\
Stock & What you hold now and what it is worth; and \textbf{Low Stock}. \\
Parties & \textbf{One} report: every party with \emph{Owes you}, \emph{You owe} and \emph{Net}. \\
Expenses & Your spending, by category and period. \\
Profit \& Loss & Whether the shop is making money over the period. \\
GST & Your tax summary (only if GST is on). \\
\bottomrule
\end{tabularx}
\end{center}
\shot{14}{14-reports.png}{Reports: pick a tab and a date range, then read, print or download.}

\section{The Party report}
Because a party can be both a customer and a supplier, the Parties report puts
them in \textbf{one list}: what each \textbf{owes you}, what \textbf{you owe} them,
and the \textbf{net}. From each row you can jump straight to \ui{Receive} or
\ui{Pay}.

\section{The Party Statement (and PDF / Excel)}
A party's \textbf{statement} is their full account on one page. It shows, in date
order:
\begin{itemize}[leftmargin=1.3em,itemsep=2pt]
  \item every \textbf{sale}, \textbf{purchase}, \textbf{receipt} and \textbf{payment};
  \item the \textbf{payment method} shown beneath each receipt or payment;
  \item a running \textbf{receivable} balance and a running \textbf{payable} balance, side by side;
  \item each document's \textbf{outstanding} amount.
\end{itemize}
Download it as a \textbf{PDF} (to print or send on WhatsApp) or \textbf{Excel}
(to work with the numbers). Both are produced ready-to-share.
\shot{15}{15-party-statement.png}{A party statement with running receivable and payable balances, ready to download.}

% =====================================================================
\chapter{Your Daily Shop Routine}
% =====================================================================
A simple rhythm keeps your books correct with almost no effort.

\begin{center}
\small
\begin{tabularx}{\linewidth}{>{\bfseries\color{brandgreendark}}l X}
\toprule
Morning & Sign in $\rightarrow$ glance at the dashboard $\rightarrow$ check \textbf{Low Stock} $\rightarrow$ check who still owes you. \\
\addlinespace
During the day & Record \textbf{every} sale and purchase as it happens; record money \textbf{received} and \textbf{paid}; adjust stock only when something is broken, lost or counted. \\
\addlinespace
At closing & Review today's \textbf{sales}, \textbf{cash received} and \textbf{money paid}; record the day's \textbf{expenses}; check your \textbf{Cash \& Bank} balance; look at outstanding parties. \\
\bottomrule
\end{tabularx}
\end{center}

\begin{tipbox}
The golden rule: \textbf{record it when it happens}. A sale entered at the moment
of billing is a sale you never have to remember later.
\end{tipbox}

% =====================================================================
\chapter{Common Mistakes and Fixes}
% =====================================================================
\begin{center}
\small
\begin{tabularx}{\linewidth}{>{\bfseries\color{navy}}p{5.2cm} X}
\toprule
Problem & What to do \\
\midrule
Product not showing in billing & Make sure it is \textbf{added} and \textbf{active} in Products / Items. \\
Wrong quantity or price on a draft & Edit the draft before you complete it. \\
Wrong price already completed & A completed invoice cannot be edited — use a \textbf{sales return}, or correct with a new entry. \\
Wrong party selected & On a draft, change it; if already posted, reverse with a return and re-enter. \\
Payment not reducing a bill & Check the amount — it fills the \textbf{oldest open bills} first; a smaller amount only partly pays. \\
``Looks like a bill already recorded'' & You are about to enter the \textbf{same} bill twice. Open the existing one, or confirm again only if it really is new. \\
Bill reading failed & Retry the upload, or enter the bill with the normal Purchase / Sales form. \\
Stock looks wrong & Use \textbf{Stock Adjustments} after a physical count to set it right. \\
A report looks empty & Check the \textbf{date range} at the top of the report. \\
Cannot download PDF / Excel & Allow pop-ups / downloads for the site in your browser, then try again. \\
Signed out unexpectedly & Sign in again; your data is safe. \\
\bottomrule
\end{tabularx}
\end{center}

\begin{notebox}
Completed (posted) documents are kept unchangeable on purpose — it is what keeps
your books trustworthy. The correct way to fix them is a \textbf{return} or an
\textbf{adjustment}, which leaves a clean trail.
\end{notebox}

If something still looks wrong, contact your Byapar support contact (see the Quick
Reference).

% =====================================================================
\chapter{Frequently Asked Questions}
% =====================================================================
\begin{faqbox}{Can one party both buy from me and sell to me?}
Yes. A party is one record that holds both sides. Byapar keeps ``owes you'' and
``you owe'' separately and never cancels one against the other.
\end{faqbox}
\begin{faqbox}{Can I sell on credit?}
Yes. Choose \ui{On credit} (or \ui{Part}) on the bill. The unpaid amount becomes
what the customer owes you, and appears in the Credit Book.
\end{faqbox}
\begin{faqbox}{Can I record a part payment?}
Yes. Enter whatever was paid in Money Received or Money Paid; it settles the
oldest bills first, and the rest stays open.
\end{faqbox}
\begin{faqbox}{Can I upload a handwritten bill?}
Yes, if bill reading is switched on for your shop. Neat handwriting reads best —
always check the figures before you confirm.
\end{faqbox}
\begin{faqbox}{Can I download statements and reports?}
Yes. Party statements and most reports download as PDF or Excel.
\end{faqbox}
\begin{faqbox}{Can I see what is running low?}
Yes. The \textbf{Low Stock} page lists every item at or below its minimum level.
\end{faqbox}
\begin{faqbox}{Can I use Byapar without GST?}
Yes. Leave GST off and your bills carry no tax — just the item amounts.
\end{faqbox}
\begin{faqbox}{What happens when I make a sale?}
Stock goes down, an invoice is created, any unpaid amount becomes what the customer
owes you, and the money paid is recorded.
\end{faqbox}
\begin{faqbox}{What happens when I record a purchase?}
Stock goes up at its cost, and any amount you did not pay becomes what you owe the
supplier.
\end{faqbox}
\begin{faqbox}{Can I correct a transaction after completing it?}
A completed document is final. Reverse it with a \textbf{return} (sales or purchase)
or a \textbf{stock adjustment}, then re-enter it correctly.
\end{faqbox}
\begin{faqbox}{What if I forget to record a payment?}
Record it whenever you remember — the party's balance updates straight away. Enter
the real date so the statement stays accurate.
\end{faqbox}

% =====================================================================
\chapter{Quick Reference}
% =====================================================================
\section{Where things live}
\begin{center}
\small
\begin{tabularx}{\linewidth}{>{\bfseries\color{navy}}l X}
\toprule
Make a quick bill & Overview $\rightarrow$ Quick Billing \\
Full sale / invoices & Sales $\rightarrow$ New sale / Invoices \\
Record a purchase & Purchase $\rightarrow$ Purchases \\
What you owe & Purchase $\rightarrow$ Payables \\
Add / view parties & Parties $\rightarrow$ Add Party / All Parties \\
Who owes you & Parties $\rightarrow$ Credit Book \\
Take / make a payment & Money $\rightarrow$ Money Received / Money Paid \\
Spending & Money $\rightarrow$ Expenses; balances in Cash \& Bank \\
Stock \& low stock & Stock \& Inventory $\rightarrow$ Stock Dashboard / Low Stock \\
Scan a paper bill & Overview $\rightarrow$ Bill Import \\
Reports \& statements & Reports (Sales, Parties, Stock, Profit \& Loss\ldots) \\
\bottomrule
\end{tabularx}
\end{center}

\section{Daily checklist}
\textbf{Morning:} sign in; dashboard; Low Stock; who owes you.\\
\textbf{Through the day:} record every sale, purchase, receipt and payment.\\
\textbf{Closing:} review sales, cash and payments; record expenses; check Cash \&
Bank; review outstanding parties.

\section{Words to know}
\begin{center}
\small
\begin{tabularx}{\linewidth}{>{\bfseries\color{navy}}l X}
\toprule
Party & Any person or business you deal with (customer, supplier, or both). \\
Receivable / Owes you & Money a customer still has to pay you. \\
Payable / You owe & Money you still have to pay a supplier. \\
Stock & How much of a product you have right now. \\
Draft & A saved entry that changes nothing until you complete it. \\
Completed / Posted & A final entry — stock and money have moved. \\
Advance & Money received beyond the open bills, kept for next time. \\
\bottomrule
\end{tabularx}
\end{center}

\section{Support}
For help, contact your Byapar support contact. \emph{(Your shop's support email and
phone are set by your business; if they are shown on your landing page, use those.)}

\vfill
\begin{center}
{\color{brandgreen}\rule{0.5\linewidth}{1pt}}\\[6pt]
{\sffamily\color{inkgray}\small Thank you for running your shop with \textbf{\color{navy}Byapar}.}
\end{center}

\end{document}
